\subsection{Variadic Positional Parameters}

A normal positional parameter receives one argument object. A variadic positional parameter receives any remaining positional argument objects and packs them into one tuple.

The syntax uses a single star in the function definition:

\begin{verbatim}
def func(*args):
    ...
\end{verbatim}

The name \texttt{args} is only a convention. The star is the important syntax. Python programmers usually use \texttt{args} because it clearly means ``positional arguments collected here.''

This mechanism does not create a special argument object. It creates a normal tuple.

\subsubsection{Argument Packing Mechanics: Collecting Positional Overflow into \texttt{*args}}

A function can collect an arbitrary number of positional arguments using \texttt{*args}.

\begin{verbatim}
def show_args(*args):
    print(f"args:       {args}")
    print(f"type(args): {type(args)}")

show_args("Homer", "Bart", 42)
\end{verbatim}

Possible output:

\begin{verbatim}
args:       ('Homer', 'Bart', 42)
type(args): <class 'tuple'>
\end{verbatim}

The call supplied three positional arguments. The function definition had one variadic positional parameter. Python packed the three argument objects into one tuple and bound the local name \texttt{args} to that tuple.

The binding shape is:

\begin{verbatim}
show_args("Homer", "Bart", 42)

inside function frame:
    args -> ("Homer", "Bart", 42)
\end{verbatim}

The tuple is a normal tuple object.

\begin{verbatim}
def show_args(*args):
    print(f"type(args): {type(args)}")
    print()

    selected_names = [
        "count",
        "index",
        "__len__",
        "__getitem__",
        "__contains__",
    ]

    for name in selected_names:
        print(f"{name:<14} {name in dir(args)}")

    print()
    print(f"len(args):        {len(args)}")
    print(f"args[0]:          {args[0]!r}")
    print(f"args.count(42):   {args.count(42)}")
    print(f"42 in args:       {42 in args}")

show_args("Homer", "Bart", 42)
\end{verbatim}

Possible output:

\begin{verbatim}
type(args): <class 'tuple'>

count          True
index          True
__len__        True
__getitem__    True
__contains__   True

len(args):        3
args[0]:          'Homer'
args.count(42):   1
42 in args:       True
\end{verbatim}

The tuple is immutable. The function cannot append directly to \texttt{args} because tuples do not have \texttt{append()}.

\begin{verbatim}
def show_args(*args):
    print(f"args: {args}")
    print(f"'append' in dir(args): {'append' in dir(args)}")

    try:
        args.append("D'oh!")
    except AttributeError as err:
        print(f"error type: {type(err)}")
        print(f"message:    {err}")

show_args("Homer", 42)
\end{verbatim}

Possible output:

\begin{verbatim}
args: ('Homer', 42)
'append' in dir(args): False
error type: <class 'AttributeError'>
message:    'tuple' object has no attribute 'append'
\end{verbatim}

The tuple itself cannot be changed in place, but the objects inside the tuple may still be mutable. The tuple stores references to the argument objects.

\begin{verbatim}
def add_quote(*args):
    print(f"args:       {args}")
    print(f"type(args): {type(args)}")

    first_item = args[0]
    print(f"type(first_item): {type(first_item)}")

    first_item.append("Nobody expects the Spanish Inquisition!")

quotes = ["D'oh!", "No soup for you!"]

print("before call")
print(f"quotes:     {quotes}")
print(f"id(quotes): {id(quotes)}")
print()

add_quote(quotes)

print()
print("after call")
print(f"quotes:     {quotes}")
print(f"id(quotes): {id(quotes)}")
\end{verbatim}

Possible output:

\begin{verbatim}
before call
quotes:     ["D'oh!", 'No soup for you!']
id(quotes): 2276834282880

args:       (["D'oh!", 'No soup for you!'],)
type(args): <class 'tuple'>
type(first_item): <class 'list'>

after call
quotes:     ["D'oh!", 'No soup for you!', 'Nobody expects the Spanish Inquisition!']
id(quotes): 2276834282880
\end{verbatim}

The tuple \texttt{args} was not mutated. The list object inside the tuple was mutated. This follows the same rule from the previous section: function calls share object references. The \texttt{*args} tuple is only the container used to collect the remaining positional references.

A variadic positional parameter can appear after ordinary positional parameters.

\begin{verbatim}
def show_report(title, *lines):
    print(f"title: {title}")
    print(f"lines: {lines}")
    print(f"type(lines): {type(lines)}")

show_report(
    "Important facts",
    "The answer is 42.",
    "Chuck Norris can divide by zero.",
    "It is only a model.",
)
\end{verbatim}

Possible output:

\begin{verbatim}
title: Important facts
lines: ('The answer is 42.', 'Chuck Norris can divide by zero.', 'It is only a model.')
type(lines): <class 'tuple'>
\end{verbatim}

The first positional argument is bound normally:

\begin{verbatim}
title -> "Important facts"
\end{verbatim}

The remaining positional arguments are packed into the variadic tuple:

\begin{verbatim}
lines -> (
    "The answer is 42.",
    "Chuck Norris can divide by zero.",
    "It is only a model."
)
\end{verbatim}

If there are no remaining positional arguments, Python creates an empty tuple.

\begin{verbatim}
def show_report(title, *lines):
    print(f"title:      {title!r}")
    print(f"lines:      {lines}")
    print(f"type(lines): {type(lines)}")
    print(f"len(lines):  {len(lines)}")

show_report("Empty report")
\end{verbatim}

Possible output:

\begin{verbatim}
title:      'Empty report'
lines:      ()
type(lines): <class 'tuple'>
len(lines):  0
\end{verbatim}

An empty \texttt{*args} value is not \texttt{None}. It is an empty tuple.

The parameter name does not have to be \texttt{args}.

\begin{verbatim}
def show_people(*people):
    print(f"people:       {people}")
    print(f"type(people): {type(people)}")

show_people("Jerry", "George", "Elaine", "Kramer")
\end{verbatim}

Possible output:

\begin{verbatim}
people:       ('Jerry', 'George', 'Elaine', 'Kramer')
type(people): <class 'tuple'>
\end{verbatim}

The star makes the parameter variadic. The name only chooses the local name inside the function frame.

The standard \texttt{inspect} module shows this parameter kind explicitly.

\begin{verbatim}
import inspect

def show_report(title, *lines):
    print(title, lines)

sig = inspect.signature(show_report)

print(f"signature: {sig}")
print()

for name, param in sig.parameters.items():
    print(f"{name:<8} kind={param.kind}")
\end{verbatim}

Possible output:

\begin{verbatim}
signature: (title, *lines)

title    kind=POSITIONAL_OR_KEYWORD
lines    kind=VAR_POSITIONAL
\end{verbatim}

The \texttt{VAR\_POSITIONAL} kind means that the parameter collects positional overflow.

A common use is to process all remaining values in a loop.

\begin{verbatim}
def print_total(label, *values):
    total = 0

    print(f"label: {label}")

    for value in values:
        print(f"adding {value}")
        total = total + value

    print(f"total: {total}")

print_total("scores", 10, 20, 12)
\end{verbatim}

Possible output:

\begin{verbatim}
label: scores
adding 10
adding 20
adding 12
total: 42
\end{verbatim}

Here, \texttt{values} is the tuple \texttt{(10, 20, 12)}. The loop iterates over that tuple.

This is still ordinary tuple iteration. There is no hidden magic after the packing step.

\subsubsection{Argument Unpacking Operations: Expanding Sequences Across Function Call Boundaries with \texttt{*}}

The single star has a second use at the call site. In a function definition, \texttt{*args} packs positional overflow into a tuple. In a function call, \texttt{*some\_iterable} unpacks an iterable object into separate positional arguments.

These are opposite directions:

\begin{verbatim}
definition side:
    def func(*args):
        pack many positional arguments into one tuple

call side:
    func(*items)
        unpack one iterable into many positional arguments
\end{verbatim}

Start with a normal function.

\begin{verbatim}
def show_line(name, score, mood):
    print(f"name={name!r}")
    print(f"score={score!r}")
    print(f"mood={mood!r}")

data = ["Homer", 42, "hungry"]

show_line(*data)
\end{verbatim}

Possible output:

\begin{verbatim}
name='Homer'
score=42
mood='hungry'
\end{verbatim}

The call:

\begin{verbatim}
show_line(*data)
\end{verbatim}

means:

\begin{verbatim}
take elements from data
send them as separate positional arguments
\end{verbatim}

So this:

\begin{verbatim}
show_line(*data)
\end{verbatim}

is equivalent to:

\begin{verbatim}
show_line(data[0], data[1], data[2])
\end{verbatim}

for this specific three-element list.

The list object can be inspected.

\begin{verbatim}
data = ["Homer", 42, "hungry"]

print(f"type(data): {type(data)}")
print()

selected_names = [
    "append",
    "extend",
    "__len__",
    "__getitem__",
    "__iter__",
]

for name in selected_names:
    print(f"{name:<14} {name in dir(data)}")

print()
print(f"len(data):       {len(data)}")
print(f"data[0]:         {data[0]!r}")
print(f"list iteration:  {[item for item in data]}")
\end{verbatim}

Possible output:

\begin{verbatim}
type(data): <class 'list'>

append         True
extend         True
__len__        True
__getitem__    True
__iter__       True

len(data):       3
data[0]:         'Homer'
list iteration:  ['Homer', 42, 'hungry']
\end{verbatim}

The unpacking operation uses iteration. The object must be iterable.

A tuple can also be unpacked.

\begin{verbatim}
def show_line(name, score, mood):
    print(f"{name:<10} score={score:04d} mood={mood}")

data = ("Lisa", 100, "focused")

print(f"type(data): {type(data)}")
show_line(*data)
\end{verbatim}

Possible output:

\begin{verbatim}
type(data): <class 'tuple'>
Lisa       score=0100 mood=focused
\end{verbatim}

A range can also be unpacked because it is iterable.

\begin{verbatim}
def show_three(a, b, c):
    print(f"a={a}, b={b}, c={c}")

nums = range(1, 4)

print(f"nums:       {nums}")
print(f"type(nums): {type(nums)}")
print()

selected_names = [
    "__len__",
    "__getitem__",
    "__iter__",
    "count",
    "index",
]

for name in selected_names:
    print(f"{name:<14} {name in dir(nums)}")

print()
show_three(*nums)
\end{verbatim}

Possible output:

\begin{verbatim}
nums:       range(1, 4)
type(nums): <class 'range'>

__len__        True
__getitem__    True
__iter__       True
count          True
index          True

a=1, b=2, c=3
\end{verbatim}

The \texttt{range} object is not a list, but it can produce values one by one. The star operator can expand those produced values into positional arguments.

A string can also be unpacked. This is sometimes useful, but often surprising.

\begin{verbatim}
def show_three(a, b, c):
    print(f"a={a!r}")
    print(f"b={b!r}")
    print(f"c={c!r}")

text = "Ni!"

print(f"type(text): {type(text)}")
show_three(*text)
\end{verbatim}

Possible output:

\begin{verbatim}
type(text): <class 'str'>
a='N'
b='i'
c='!'
\end{verbatim}

The string is iterable over characters, so \texttt{*text} expands it into separate character arguments.

The number of unpacked items must match the receiving signature unless the receiving function has a variadic positional parameter.

\begin{verbatim}
def show_three(a, b, c):
    print(f"a={a}, b={b}, c={c}")

short_data = [1, 2]
long_data = [1, 2, 3, 4]

try:
    show_three(*short_data)
except TypeError as err:
    print("short_data failed")
    print(f"error type: {type(err)}")
    print(f"message:    {err}")
    print()

try:
    show_three(*long_data)
except TypeError as err:
    print("long_data failed")
    print(f"error type: {type(err)}")
    print(f"message:    {err}")
\end{verbatim}

Possible output:

\begin{verbatim}
short_data failed
error type: <class 'TypeError'>
message:    show_three() missing 1 required positional argument: 'c'

long_data failed
error type: <class 'TypeError'>
message:    show_three() takes 3 positional arguments but 4 were given
\end{verbatim}

The unpacking operation expands the iterable first. Then normal argument binding rules apply.

A variadic receiving function can accept any number of unpacked positional arguments.

\begin{verbatim}
def show_all(*items):
    print(f"items:       {items}")
    print(f"type(items): {type(items)}")
    print(f"len(items):  {len(items)}")

data = ["Arthur", "Patsy", "Robin", "Galahad"]

show_all(*data)
\end{verbatim}

Possible output:

\begin{verbatim}
items:       ('Arthur', 'Patsy', 'Robin', 'Galahad')
type(items): <class 'tuple'>
len(items):  4
\end{verbatim}

This call performs two opposite operations in sequence:

\begin{verbatim}
show_all(*data)

call side:
    unpack data into separate positional arguments

definition side:
    pack those positional arguments into the local tuple items
\end{verbatim}

For this example, the overall effect looks similar to converting the list into a tuple. But conceptually, two separate call-boundary operations happened.

Unpacking does not copy the contained objects. It passes references to the same objects.

\begin{verbatim}
def mutate_first(*items):
    first = items[0]
    first.append("It is only a model.")

quotes = ["D'oh!"]
data = [quotes]

print("before")
print(f"quotes:     {quotes}")
print(f"id(quotes): {id(quotes)}")
print()

mutate_first(*data)

print("after")
print(f"quotes:     {quotes}")
print(f"id(quotes): {id(quotes)}")
\end{verbatim}

Possible output:

\begin{verbatim}
before
quotes:     ["D'oh!"]
id(quotes): 2276834282880

after
quotes:     ["D'oh!", 'It is only a model.']
id(quotes): 2276834282880
\end{verbatim}

The list \texttt{data} was unpacked. Its first element was the list object also referenced by \texttt{quotes}. The function received that same inner list object and mutated it.

The star can also be mixed with normal positional arguments.

\begin{verbatim}
def show_line(name, score, mood):
    print(f"name={name!r}")
    print(f"score={score!r}")
    print(f"mood={mood!r}")

tail = [42, "hungry"]

show_line("Homer", *tail)
\end{verbatim}

Possible output:

\begin{verbatim}
name='Homer'
score=42
mood='hungry'
\end{verbatim}

The call expands to:

\begin{verbatim}
show_line("Homer", 42, "hungry")
\end{verbatim}

Multiple starred iterables can appear in one call.

\begin{verbatim}
def show_all(*items):
    for index, item in enumerate(items, start=1):
        print(f"{index:>2}. {item}")

first = ["Homer", "Bart"]
second = ["Lisa", "Marge"]
third = ["Maggie", 42]

show_all(*first, *second, *third)
\end{verbatim}

Possible output:

\begin{verbatim}
 1. Homer
 2. Bart
 3. Lisa
 4. Marge
 5. Maggie
 6. 42
\end{verbatim}

Python expands each starred iterable into positional arguments from left to right.

The argument stream seen by \texttt{show\_all()} is:

\begin{verbatim}
"Homer", "Bart", "Lisa", "Marge", "Maggie", 42
\end{verbatim}

The receiving \texttt{*items} parameter then packs that stream into one tuple.

A non-iterable object cannot be unpacked with \texttt{*}.

\begin{verbatim}
def show_all(*items):
    print(items)

value = 42

try:
    show_all(*value)
except TypeError as err:
    print(f"error type: {type(err)}")
    print(f"message:    {err}")
\end{verbatim}

Possible output:

\begin{verbatim}
error type: <class 'TypeError'>
message:    __main__.show_all() argument after * must be an iterable, not int
\end{verbatim}

The exact module name in the error message may differ. The rule is stable: call-side \texttt{*} requires an iterable object.

The compact model is:

\begin{verbatim}
def f(*args):
    star in definition packs positional overflow into a tuple

f(*items):
    star in call unpacks an iterable into positional arguments
\end{verbatim}

The two uses are mirror images. One collects many arguments into one tuple. The other expands one iterable into many arguments.